\section{Approach}
\label{sec:approach}

\subsection{Overview}

Let $N$ be an NTA and
$\Phi=\{\varphi_1,\ldots,\varphi_m\}$ a set of timed safety
properties. We write $M=(N,\Phi)$ for the repair input.
MPTA-Repair changes only repairable numerical clock bounds in guards
and invariants. It returns a repaired specification
$M'=(N',\Phi)$ only when $N'$ satisfies the accepted-repair condition
in Eq.~\eqref{eq:accepted-repair}; otherwise, it reports that no repair
was found within the resource budget.

The main difficulty is that the properties in $\Phi$ do not define
independent repair tasks. Different properties may constrain the same
clock bound, overlapping execution paths, or different stages of a
shared timing budget. Their TDTs may also have different relations:
some are duplicates, some share only part of a discrete path, and some
reach different violations through the same repairable bounds.
Consequently, an edit derived from one TDT may leave another violating
execution feasible or invalidate a property that already holds. A
single-property workflow that processes these TDTs independently does
not retain their joint constraints and may repeatedly explore
incompatible candidate edits.

MPTA-Repair uses these relations to organize the current repair
problem before MaxSMT solving. At the property level, equivalent or
stronger requirements are identified to avoid redundant constraints.
At the trace level, duplicate and structurally related TDTs are grouped.
At the repair-bound level, each property-trace constraint is mapped to
the bounds that can affect it. Property-trace constraints that depend
on disjoint bound sets form separate search groups, whereas constraints
that share repairable bounds remain in the same group and are solved
jointly. This partition reduces independent parts of the search without
separating constraints that require coordinated edits.

A trace-local repair workflow encodes one TDT, solves a MaxSMT
instance, and validates the resulting candidate. MPTA-Repair extends
this workflow with a shared property-trace pool, relation-guided
selection of active constraints and bounds, and iterative refinement
across properties. At refinement round $k$,
$\mathcal{T}_k$ denotes the current set of property-trace pairs
$(\tau,\varphi)$, where $\tau$ is a TDT for
$\varphi\in\Phi$. The framework verifies $N'$ against the full property
set, updates $\mathcal{T}_k$, constructs a joint MaxSMT problem for the
dependent groups, and instantiates the resulting clock-bound edits in a
candidate model.

A candidate is not accepted merely because it satisfies the active
MaxSMT constraints. It must first pass full-property regression
verification and then the admissibility check based on functional
equivalence. If either check fails, newly exposed TDTs are added to the
pool and the rejected assignment is excluded from later rounds. Only a
candidate that passes both checks is returned as $M'$.
Figure~\ref{fig:mpta-repair-workflow} summarizes this iterative
workflow.

\begin{figure*}[t]
  \centering
  \includegraphics[width=0.86\textwidth]{MPTA_Repair_overview.pdf}
  \caption{Overview of the MPTA-Repair workflow.}
  \label{fig:mpta-repair-workflow}
\end{figure*}

\subsection{Joint Trace Encoding and Optimization Objective}

MPTA-Repair follows the clock-bound variation strategy used in TARTAR
\cite{kolbl2019clock}. Let $S$ denote the set of repairable clock-bound
occurrences. For each occurrence $s\in S$, the original numerical bound
$b_s$ in a transition guard or location invariant is associated with a
variation variable $\delta_s$. This ensures that the repaired constraint
retains the original clock reference and comparison operator while
incorporating a shifted value:
\begin{equation}
  b'_s = b_s + \delta_s
\end{equation}
The repair space is therefore strictly restricted to clock-bound
constants, keeping transition structures and discrete updates unaltered.

MPTA-Repair builds on the trace-encoding formulation of TARTAR and
extends it to the multi-property setting \cite{kolbl2019clock}. Given
a TDT
\begin{equation}
\begin{aligned}
  \tau ={}&
  \ell_0 \xrightarrow{d_0,e_0} \ell_1 \cdots \\
  &{}\xrightarrow{d_{n-1},e_{n-1}} \ell_n
\end{aligned}
\end{equation}
a trace path is encoded as a conjunction of linear arithmetic
constraints:
\begin{equation}
\begin{aligned}
  \operatorname{Path}(\tau,\rho,\mathbf{z}_{\tau})
  ={}& \mathit{Init}(\nu_0)\\
  &{}\land
  \bigwedge_{i=0}^{n-1}\Bigl(
    d_i \ge 0
    \land \mathit{Inv}(\ell_i,\nu_i+d_i,\rho)\\
  &{}\land
    \mathit{Guard}(e_i,\nu_i+d_i,\rho)\\
  &{}\land
    \nu_{i+1}
    =
    \mathit{Reset}(\nu_i+d_i,R_i)
  \Bigr)
\end{aligned}
\end{equation}
Here, $\rho$ denotes the modified bounds, and $\mathbf{z}_{\tau}$
collects the delays and clock valuations. The reset equation
$\nu_{i+1}=\mathit{Reset}(\nu_i+d_i,R_i)$ updates the valuation after
transition $e_i$ by resetting clocks in $R_i$ and preserving the others.
The original violation encoding combines path executability with the
final property violation:
\begin{equation}
  \operatorname{Enc}(\tau,\varphi,\rho)
  =
  \operatorname{Path}(\tau,\rho,\mathbf{z}_{\tau})
  \land
  \mathit{Bad}_{\varphi}(\ell_n,\nu_n)
\end{equation}
The predicate $\mathit{Bad}_{\varphi}(\ell_n,\nu_n)$ characterizes the
violation state of $\varphi$ at the final location. To block this
violating behavior while preserving executability of the discrete path
skeleton, the single-trace repair constraint uses two copies of the
trace variables:
\begin{equation}
\begin{aligned}
  &\Psi_{\tau,\varphi}^{\mathrm{repair}}(\rho)
  =
  \exists \mathbf{z}_{\tau}^{+}.
  \operatorname{Path}(\tau,\rho,\mathbf{z}_{\tau}^{+}) \\
  &\quad{}\land
  \neg \exists \mathbf{z}_{\tau}^{-}.
  \Bigl(
    \operatorname{Path}(\tau,\rho,\mathbf{z}_{\tau}^{-}) \\
  &\qquad{}\land
    \mathit{Bad}_{\varphi}(\ell_n^{-},\nu_n^{-})
  \Bigr)
\end{aligned}
\end{equation}
The first part requires that the original discrete path skeleton still
has at least one timed realization after repair. The second part
requires that no assignment of delays and clock valuations can make that
same path violate $\varphi$.

We extend this single-trace formulation to a joint, multi-property
incremental setting. At refinement round $k$, counterexamples from all
currently violated properties are accumulated into a trace pool
\begin{equation}
  \mathcal{T}_k=\{(\tau_j,\varphi_j)\}_{j=1}^{m_k}
\end{equation}
MPTA-Repair assembles a joint MaxSMT problem with the constraint system
defined as
\begin{equation}
\begin{aligned}
  \mathcal{F}_k(\rho)
  =
  &\mathit{Dom}(\rho)
  \land \mathit{Block}_k(\rho) \\
  &{}\land
  \bigwedge_{(\tau,\varphi)\in \mathit{Select}(\mathcal{T}_k)}
  \Psi_{\tau,\varphi}^{\mathrm{repair}}(\rho)
\end{aligned}
\end{equation}
Here, $\mathit{Dom}(\rho)$ enforces static interval consistency and
domain limits, such as $b'_s\ge 0$, and
$\mathit{Select}(\mathcal{T}_k)$ represents a reduced subset of traces
obtained after relation-aware filtering. When an instantiated candidate
NTA fails verification, the constraint system is updated
incrementally. The framework expands the pool via
\begin{equation}
  \mathcal{T}_{k+1}
  =
  \mathcal{T}_k
  \cup
  \mathit{NewTDTs}(N(\rho_k),\Phi)
\end{equation}
and updates the search-space block to exclude the failed assignment
$\rho_k$ over the set of repairable bounds $S$:
\begin{equation}
  \mathit{Block}_{k+1}(\rho)
  =
  \mathit{Block}_k(\rho)
  \land
  \left(
    \bigvee_{s\in S} \delta_s \ne \delta_s^k
  \right)
\end{equation}
where $\delta_s^k$ is the valuation assigned to variation variable
$\delta_s$ in round $k$.

The optimization objective is governed by a hierarchical set of soft
constraints. MPTA-Repair uses a lexicographic objective that first
favors small repairs and then incorporates risk- and coverage-based
priorities:
\begin{equation}
\operatorname*{lex\,min}_{\rho}
\left(
\begin{aligned}
  &\sum_{s\in S}\mathbf{1}\{\delta_s \ne 0\},\\
  &\sum_{s\in S}|\delta_s|,\\
  &\sum_{s\in S} \mathit{risk}_s \cdot \mathbf{1}\{\delta_s \ne 0\},\\
  &-\sum_{s\in S} \mathit{benefit}_s \cdot \mathbf{1}\{\delta_s \ne 0\}
\end{aligned}
\right)
\end{equation}
Here, $\mathbf{1}\{\cdot\}$ denotes the indicator function, yielding 1
when the inner condition holds and 0 otherwise. The primary and
secondary objectives minimize the number and total magnitude of altered
clock bounds to produce concise repairs. The third term incorporates a
penalty coefficient $\mathit{risk}_s$ to discourage modifications prone
to regression, while the final term uses a reward coefficient
$\mathit{benefit}_s$ to prioritize repair sites that cover a broader set
of violated traces.

\subsection{Relation-Guided Search Optimization}

During refinement, accumulating violated properties, TDTs, and editable
bounds expands the MaxSMT search space. To reduce redundant
constraints, MPTA-Repair uses relation-guided optimization to filter,
group, and rank constraints and repair variables. These heuristics only
guide the search; correctness is still guaranteed through full-property
regression verification and an admissibility check.

At the property layer, timed safety properties are normalized into an
$\mathtt{A}[]$ fragment over location predicates and clock or integer
constraints. Equivalence is detected by the syntactic identity of
normalized forms, while dominance is evaluated under logical
subsumption, e.g., $\mathtt{A}[]\, x \le 5$ subsumes
$\mathtt{A}[]\, x \le 10$ within the same structural context.
MPTA-Repair then retains only the stronger properties for solving to
optimize efficiency.

At the counterexample layer, MPTA-Repair maintains a TDT pool indexed by
violated properties. Exact duplicates are removed, and the remaining
traces are annotated with traversed transitions, visited locations,
active clocks, and guard or invariant bounds. Traces with identical
discrete transition sequences or overlapping clock zones are grouped to
reveal common bottlenecks. A trace constraint is omitted from the active
encoding only when it is an exact duplicate or soundly subsumed by an
encoded trace constraint.

At the repair-bound layer, MPTA-Repair maps each active property-trace
constraint to the clock bounds occurring in or affecting its encoding.
This mapping restricts active MaxSMT decision variables to bounds
relevant to the current counterexample pool and ranks them for candidate
generation. Bounds with higher relational scores are prioritized
because their modifications are more likely to satisfy the joint
constraint system.